\section{Discussion}

This study treats molecular benchmark construction as part of the measurement process. The primary intervention was not a generic comparison of unrelated splitters: within a fixed target-blind candidate pool, a scaffold-disjoint partition with lower train--test target-mean mismatch was selected only when it had exactly the same test-set size as its target-blind size-matched counterpart. Classification provided no corrected evidence of a target-balanced advantage, whereas primary regression effects were large under one acyclic-scaffold convention and weakened sharply under another. Molecular-record perturbation preserved the corrected classification conclusion while reversing most point-estimate directions. Together, these results show that benchmark effects are protocol-conditioned rather than model-intrinsic.

\subsection{Target-mean-aware selection is an intervention on the benchmark, not an isolated causal variable}

The paired design controls dataset, partition seed, scaffold rule, candidate-generation budget, and test cardinality, but it does not hold the selected molecular subset fixed. A different candidate can therefore carry different scaffold concentration, acyclic content, response-shape diagnostics, and other structural properties. The collateral diagnostics make this explicit. Accordingly, the estimand is the performance consequence of a target-mean-aware partition-selection rule under controlled search conditions, not the causal effect of ``target distribution'' in isolation.

This distinction matters for interpretation. A lower target-mean gap can make a test population statistically closer to the training population in one dimension while simultaneously moving the benchmark to another part of chemical space. A target-balanced split should therefore be treated as an audit perturbation rather than a universally preferable leaderboard split. Temporal or simulated-time designs remain more appropriate when prospective medicinal-chemistry deployment is the scientific target, while leakage-reduced approaches such as DataSAIL address a different validation objective \citep{landrum2023simpd,joeres2025datasail}.

\subsection{The regression effect contains a benchmark-geometry component}

The mean-only control clarifies why regression requires special care. A predictor that uses no molecular features and always returns the training-set mean has test MSE equal to the test variance plus the squared train--test mean difference. Because the partition-selection rule directly minimizes that mean difference, part of any RMSE improvement can arise mechanically from the geometry of the response distribution. The learned-model regression effects therefore cannot be interpreted purely as evidence that molecular structure generalizes better under target-mean balancing.

This finding strengthens rather than weakens the chemometric message. It shows that benchmark construction can change apparent predictive difficulty before a molecular model contributes any additional information. The appropriate comparison is consequently not simply ``balanced split versus ordinary split'', but the learned-model effect relative to a trivial response-only control and to the structural properties of the selected test population.

\subsection{Search budget and scaffold semantics are benchmark hyperparameters}

An optimized split is defined not only by its objective but also by how hard the algorithm searches. ESOL and FreeSolv continued to yield lower selected target-mean gaps as the candidate pool expanded to 20,000 partitions. Freezing search effort before predictive outcomes therefore removes an otherwise hidden researcher degree of freedom. The particular budgets used here are not universal constants; the general requirement is that the search space, stopping rule, and pre-outcome freeze be reported.

The acyclic sensitivity exposes a second hidden hyperparameter. An empty Bemis--Murcko framework can mean one shared structural group or many molecule-specific groups, and those choices encode different notions of ``unseen scaffold''. The strong attenuation of the primary regression effects under singleton semantics shows that the phrase ``scaffold split'' is scientifically incomplete unless the handling of acyclic structures is stated explicitly.

\subsection{Molecular representation belongs to the benchmark specification}

The disconnected-component sensitivity leads to the same conclusion at the molecular-record level. Selecting a dominant fragment changed fingerprints, scaffold identities, duplicate mappings, and conflicting-label groups for a non-trivial subset of BBBP, ClinTox, and HIV. Although all nine corrected sensitivity inferences remained inconclusive, seven mean-effect directions reversed. A narrative based only on the sign of a small AUC difference could therefore change under two reasonable representation policies even when neither policy provides reliable evidence for an advantage.

This asymmetry is useful: inferential status can be robust while point estimates are representation-sensitive. Reporting both levels prevents a weak directional fluctuation from being promoted into a stable scientific claim and supports the broader view that molecular identity, endpoint aggregation, and partition definitions form part of the dataset contract \citep{nael2026dataset}.

\subsection{Implications for auditable molecular benchmarking}

A defensible molecular benchmark should record more than a split ratio and random seed. At minimum, the molecular-identity and duplicate policy, structural grouping rule, requested and realized test size, target summaries, scaffold concentration, number of unique partitions, optimization objective, candidate-search budget, and machine-readable partition identifiers should be preserved. Predeclared model-seed checks should remain robustness analyses rather than pseudo-replication. When a reasonable protocol perturbation changes the conclusion, that disagreement is evidence about the benchmark and should be reported rather than resolved by choosing the preferred result.

\section{Limitations}

The study contains six public molecular-property datasets, and only ESOL and FreeSolv are regression tasks; the 20 partition seeds quantify robustness within these fixed molecular universes rather than sampling a population of chemical datasets. Predictive models are classical Morgan-fingerprint baselines, intentionally chosen to reduce architecture-related degrees of freedom. Future work should test whether the same protocol interactions persist for learned molecular representations and additional regression endpoints.

Target-aware selection optimizes only the first moment of the response distribution and exact test size; target variance, tails, mechanistic subgroups, and structural composition are not controlled. Candidate generation is a stochastic prefix heuristic rather than exhaustive subset enumeration, and the primary regression search did not converge by the frozen 20,000-candidate cap. The single-group and singleton acyclic conventions are operational extremes rather than a complete ontology of acyclic series. The dominant-fragment mapping is likewise algorithmic and is not guaranteed to identify a chemically authoritative parent. Finally, no chronological, prospective, or external industrial test set was available, so the study diagnoses benchmark construction rather than identifying the internal split that best predicts prospective drug-discovery performance.

\section{Conclusion}

Molecular benchmark performance is conditional on choices about molecular identity, structural grouping, partition search, and statistical replication. Exact-size target-mean-aware scaffold selection did not produce corrected evidence of a classification advantage, while its apparent regression benefit depended on acyclic-scaffold semantics and included a response-geometry component visible even to a mean-only predictor. A deterministic molecular-record perturbation further changed most classification point-estimate directions without changing the corrected inferential conclusion.

The practical response is not to search for one universally superior split. Validation questions should be stated explicitly, benchmark-construction rules and search effort should be frozen before outcomes are examined, partition provenance should be preserved, trivial controls should be used when the evaluation metric couples directly to the split objective, and scientific claims should be tested against defensible protocol perturbations. Treating benchmark construction as part of chemometric measurement can reduce false confidence in molecular model comparisons and make downstream decisions more auditable.
